% =====================================================================
% Section 7 -- SYSTEM   (target: ~0.65 page)
% Compiler + independent verifier architecture; producer/consumer
% interface to the versioned governance substrate; zero-sharing
% implementation discipline and its CI enforcement (the landing evidence
% for the C5 claim -- KEEP the four CI assertions and the import-root
% disjointness verbatim).  Adapter layer, per-step compiler walk-through
% and the deviation ledger format are in the extended version.
% =====================================================================
\providecommand{\sysname}{\textsc{AsOfGov}}
\providecommand{\sem}[1]{[\![#1]\!]}

\section{SYSTEM}
\label{sec:system}

\looseness=-1\sysname{} is two programs deliberately kept apart: a \emph{binding compiler}
($\approx$3.8k Python lines, a shared core plus per-base
adapters) turning a question and a declared time point into a decision
plus a certificate, and an \emph{independent verifier} ($\approx$5.1k
lines) deciding the checking relation of \S\ref{sec:certificates} without
running, importing or consulting the compiler --- that separation is the
system contribution.

\smallskip\noindent\looseness=-1\textbf{Compiler.} Given $(\sigma,T)$ it pins
$\mathrm{ver}(T)$ with its typed commit, reads $G_v$ at that version, derives
$\alpha$, evaluates the guards in the frozen order
$\mathsf{OOV}\rightarrow\mathsf{AM}\rightarrow\mathsf{MC}$ leaving each probe
transcript in the certificate, applies the disclosure gate, and emits SQL over
the certified objects, a narrowed rewrite with its cut trace, or a typed refusal
with a witness (walk-through in~\cite{asofgov-tr}).
\emph{Registration over hard-coding}: a metric's anchor, caliber route and
coverage-domain reading must be readable from $G_v$, not a dictionary
inside the compiler (\S\ref{sec:eval-divergence} measures the cost of
violating it). \emph{No post-hoc narration}: the certificate is
emitted field by field from the decision that produced the plan.

\smallskip\noindent\looseness=-1\textbf{Verifier and the anti-certification-loop rule.} The
verifier takes the certificate envelope, the structured question, the
governance tables at the pinned version, the data and the request context, and
returns accept or reject with the failing check. The question is a shared
input to both sides --- what defeats a self-signed override --- but read for
its \emph{intent} fields only: metric, the two declared instants, scope, and,
on the 10 questions of \S\ref{sec:eval-cert}, the presented window. Every
other window coordinate is \emph{re-derived} from the
registered as-of convention rather than consumed. What
the verifier never takes is compiler internals: no templates, no shared module,
no call into the compilation entry point; every fact the compiler looked up is
\emph{re-queried} from $(G_v,D)$, so certificate fields are assertions awaiting
verification --- even neutral-looking helpers are re-implemented, since a
shared wrong month expansion cancels.

\looseness=-1 The rule is enforced mechanically. A CI check parses every verifier file and
asserts that (a)~each import resolves to the Python standard library, the
database driver, or a module inside the verifier package; (b)~no verifier file
imports any compiler location, including the pilot's per-domain compilers;
(c)~the two sides' project-internal import roots intersect in the empty set; and
(d)~no dynamic-import or path-manipulation escape hatch points at a compiler
directory. On the released artifact the intersection is empty. Under this
discipline ``the compiler is buggy but the certificate verifies'' requires
the same defect in two independent implementations
\emph{and} in the declarative checking relation; the audit surface shrinks
from a compiler to a one-page specification. A certificate travels as a JSON envelope: body
(\S\ref{sec:cert:content}), probe transcripts, machine-readable
deviation list.
